\section{Building Intuition}
\label{sec:v10-controlled-checks}

The Benoit replication is the real measurement task. Two completed
checks isolate pieces that the replication leaves entangled. Their
role is deliberately modest: they explain why the certificate and
training design should work before a completed human node audit is
available. The Markov check fixes the oracle and the sufficient state,
then shows when local supervision can replace some root supervision
under a fixed token budget. The IPW check fixes the audit estimand and
stresses the sampling design, then shows how logged propensities
recover the finite-population target while changing the variance.

\subsection{Markov Sandbox: What Local State Must Carry}
\label{sec:v10-markov-check}

The Markov sandbox keeps the tree mechanics visible. A document is a
length-\(128\) token sequence. At each position, a hidden Markov chain
selects one of \(K\) registers; the selected register emits an
observed token from a disjoint synonym palette. The learner sees the
token row, not the hidden registers. The oracle \(f^\ast\) counts
adjacent hidden-register changes.

\begin{table}[t]
\centering
\small
\begin{tabular}{@{}lcc@{}}
\toprule
 & Recoverable & Structural \\
\midrule
Hidden registers \(K\) & \(4\) & \(12\) \\
Document length & \multicolumn{2}{c}{\(128\) tokens} \\
Synonym tokens per register & \multicolumn{2}{c}{\(4\), disjoint palettes} \\
Switch probability \(h\) & \(0.039\) & \(0.079\) \\
Expected changepoints & \({\sim}5\) & \({\sim}10\) \\
Training documents & \multicolumn{2}{c}{\(10{,}240\)} \\
\bottomrule
\end{tabular}
\caption{\textbf{Markov data-generating process.} The hidden register
stays fixed with probability \(1-h\), otherwise switches to a uniformly
chosen alternative. The observed token is drawn uniformly from that
register's palette.}
\label{tab:v10-markov-dgp}
\end{table}

\begin{figure}[t]
\centering
\includegraphics[width=\linewidth]{markov_changepoint_combined.pdf}
\caption{\textbf{Counting register changes with boundary state.}
A \(12\)-token illustration with four hidden colors and synonym tokens.
The learner observes the word row. Each leaf summary records internal
count, first register, and last register. The merge adds child counts
plus the boundary correction
\(\One\{\text{left last}\ne\text{right first}\}\), detecting changes
that neither child sees alone.}
\label{fig:v10-markov-changepoint-main}
\end{figure}

The sufficient state for any span is
\[
  (\mathrm{count},\mathrm{first},\mathrm{last}),
\]
where the count records transitions inside the span and the boundary
fields let parent merges add the one transition that may cross the
child boundary:
\[
(c_L,a_L,b_L)\otimes(c_R,a_R,b_R)
=
(c_L+c_R+\One\{b_L\ne a_R\},a_L,b_R).
\]
The state is ordered and associative. A scalar child count lacks the
boundary information needed to decide whether the right edge of the
left child matches the left edge of the right child. This is the core
intuition behind local state: the node label is useful because it
supervises a mergeable representation rather than a cheaper scalar copy
of the root score.

The completed allocation grid asks the budget question directly. It
uses leaf sizes \(128,64,32,16,8\), root-label shares from \(100\%\)
down to \(10\%\), and three local-allocation policies where the tree
has non-root nodes to distinguish them. A full-document label on a
\(128\)-token sequence costs \(128\) token observations. A label on a
\(16\)-token leaf costs \(16\) token observations, so eight such leaf
labels spend the same full-document-equivalent mass as one root label.
The budget split asks whether reallocating missing root-label mass to
local count labels keeps root MAE near the all-root reference.

\begin{figure}[t]
    \centering
    \includegraphics[width=0.94\linewidth]{markov_budget_split.png}
    \caption{\textbf{Splitting a fixed token budget between root and
    local labels.} The green star is the all-root reference. The green
    line spends only the reduced root-label budget, so total budget
    shrinks. The blue, purple, and rose curves keep total supervision
    mass fixed by assigning the missing mass to local labels under
    three allocation policies.}
    \label{fig:v10-markov-budget-main}
\end{figure}

Figure~\ref{fig:v10-markov-budget-main} gives the main intuition. The
local-allocation curves stay close to the all-root reference through a
broad interior band, while the root-only curve degrades as total
supervision shrinks. The effect is strongest in the recoverable scope
and remains useful in the structural scope, where more registers and
more boundaries make small leaves harder. Appendix~\ref{app:v10-markov}
gives the full leaf-size grid and provenance. The lesson is narrow but
useful: when the local state is known and node labels measure the state
needed for composition, node supervision can substitute for some root
supervision at fixed budget.

\subsection{IPW Audit Stress Check}
\label{sec:v10-ipw-check}

The IPW stress ladder tests the estimator in
Equation~\eqref{eq:v8-ht-distortion}. It simulates hard finite
populations with Bernoulli or without-replacement sampling, balanced
or adversarial propensities, and document-nonseparable oracle losses.
The large run covers \(4\) hard cases, \(7\) document counts from
\(60\) to \(3{,}840\), \(64\) population seeds, and \(15{,}000\)
sampling trials per cell, for \(1{,}792\) raw rows.

The stress cases behave as the certificate predicts. Logged
propensity weighting keeps the preference-loss and violation-rate
biases close to zero even when the smallest joint propensities are
below \(0.001\) and maximum joint weights exceed \(1{,}000\). The
price is variance: effective sample size falls sharply in the
adversarial designs, and confidence intervals widen. At
\(3{,}840\) documents, the hardest boundary-adversarial case has
minimum joint propensity \(0.00091\), maximum joint weight
\(1{,}099\), and mean effective sample size about \(7{,}206\);
preference coverage remains \(1.000\), while violation coverage is
\(0.994\). Appendix~\ref{app:v10-ipw-stress} gives the compact
table and source path.

The Markov and IPW checks support the human node-audit path for the
Manifesto trees. They make the submission evidence tiered. The
Manifesto benchmark supports the root-observed application claim; the
teacher-trace gaps support a diagnostic claim; the Markov check
supports the local-supervision mechanism; and the IPW stress ladder
supports the audit-estimator component used by the certificate.
